\documentclass{article}
\usepackage{graphicx} % Required for inserting images
\usepackage{xcolor}      % color support
\usepackage{pagecolor}   % change page background
\usepackage{placeins} % Makes float barriers a thing
\usepackage{svg} % This makes .svg files work
\usepackage[utf8]{inputenc}
\usepackage[backend=biber,style=alphabetic]{biblatex} 

\usepackage{mathrsfs} %  For the Laplace transform
\usepackage{amsmath}
\usepackage{hyperref}


% Title stuff
\title{Tattoo Derivation}
\author{Spencer Halle}
\date{\today}


\begin{document}

\maketitle
\newpage

We start with Bayes' theory:

\begin{align*}
   P(A|B) =& \frac{P(B|A) P(A)}{P(B)}\\ 
\end{align*}

\begin{itemize}
    \item Where $P(A|B)$ is equal to the probability of some value A given observation B.
    \item Where $P(B|A)$ represents the probability of observation of observation B if value A where true
    \item Where $P(A)$ represents the probability of value A period
    \item Where $P(B)$ represents the probability of observation B period
\end{itemize}

If I want to guess the true value of a location on a number line, I could set $A=\theta$, where $\theta$ is the true location, while $x$ is the measured value. If we just want good looking curves, and the ability to scale things so it's asthetically pleasing on a graph, we don't necessarily need to worry about that denominator, as it's a constant and can be thrown out, but then that makes it so we can't say $=$, instead we have to change it to $\propto$.

\begin{align*}
   P(A|B) =& \frac{P(B|A) P(A)}{P(B)}\\ 
   P(\theta|x) =& \frac{P(x|\theta) P(\theta)}{P(x)}\\ 
   P(\theta|x) \propto& P(x|\theta) P(\theta)\\
\end{align*}

And from that, we can say that:
\begin{itemize}
    \item $P(\theta|x)$ represents the probability that a given a given value is theta if we observe x
    \item $P(x|\theta)$ represents the probability you would measure x if a given value were to be $\theta$
    \item $P(\theta)$ represents the probability of the true predicted value being at point $\theta$ based on our prior
\end{itemize}

From there, we can take what we know the Gaussian distribution looks like, and fill in some things:

\begin{align*} 
   P(x|\theta) \propto& \exp(-\frac{(\theta - \mu_0)^2}{2 \sigma_0^2})\\
   P(\theta) \propto& \exp(-\frac{(\theta - x)^2}{2 \sigma^2})\\
\end{align*}

And we can say of each of those variables that:

\begin{itemize}
    \item $\mu_0$ represents the prior's mean
    \item $\sigma_0$ represents the prior's standard deviation
    \item $\sigma$ represents the standard deviation of the measurement
\end{itemize}

Which yields the expression below:
\begin{align*}
    P(\theta|x) \propto& \exp(-\frac{(\theta - \mu_0)^2}{2 \sigma_0^2})
    exp(-\frac{(\theta - x)^2}{2 \sigma^2})\\
    P(\theta|x) \propto& \exp(-\frac{(\theta - \mu_0)^2}{2 \sigma_0^2}-\frac{(\theta - x)^2}{2 \sigma^2})\\
    P(\theta|x) \propto& \exp(-\frac{1}{2}(\frac{(\theta - \mu_0)^2}{\sigma_0^2}+\frac{(\theta - x)^2}{\sigma^2}))\\
\end{align*}

But now, I need to get those normalization constants so things look right when plotting, so maybe

And that, I can manipulate and plot, which allows me to make a figure for my tattoo, so long as I control all the variables

\end{document} 


